\chapter{基于去噪掩码自编码器的脑电表征学习方法}
\section{引言}


前述第三章和第四章从信号级生成建模的角度，分别沿着增强和约简两个维度构建了协同互补的技术方案。如双层级生成式建模理论框架所述，表征级建模以学习映射函数$h:\mathcal{X}\to\mathcal{H}$为核心目标，旨在从海量无标注数据中挖掘具有强泛化能力和任务无关性的隐空间表征。本章将沿着这一思路，从自监督学习的角度探索表征级生成建模，重点解决脑电信号分析中标注数据稀缺、跨域泛化能力不足以及噪声干扰严重等关键挑战，通过构建基于去噪掩码自编码器的通用表征学习框架，实现从信号级处理到表征级理解的跃升。

在表征级生成建模的实践中，EEG信号分析面临三个核心挑战。首先是标注数据的极度稀缺性。EEG信号的精确标注需要专业的神经科医生或EEG技师具备丰富的经验和专业知识\cite{teplan2002fundamentals, tasci2023epilepsy, kaiboriboon2012EEG}，这种专业要求使得高质量标注样本的获取成本高昂且规模受限。同时，EEG信号的个体差异显著，同一认知任务在不同被试间可能表现出完全不同的模式，进一步加剧了标注的复杂性和不一致性问题。其次是跨域泛化能力的不足。不同采集设备在电极布局、采样频率、滤波参数等方面存在显著差异\cite{styles2021towards}，导致基于特定数据集训练的模型难以适应新的设备环境或应用场景。最后是噪声干扰的普遍存在。EEG信号中包含眼电伪迹、肌电干扰、工频噪声、基线漂移等多种污染成分\cite{pijn1991chaos, brake2024neurophysiological}，这些干扰信号的幅度往往超过真实的神经信号，严重影响有效信息的提取。

深度学习技术为EEG分析提供了强大的工具。CNN能够捕获空间邻域模式\cite{lawhernEEGnet2018}，RNN擅长建模时序依赖关系，GNN可以处理复杂的空间拓扑结构\cite{pan2023st}，而Transformer架构则能够捕获长程时间依赖\cite{defossez2023decoding, songEEG2022, bagchi2022eeg}。这些方法主要依赖监督学习范式\cite{mahmood2019fully}，在面对标注数据稀缺和噪声干扰严重的现实挑战时表现有限。特别是在BCI应用\cite{chaudhary2016brain}和临床诊断等对可靠性要求高的场景中，这些方法的泛化能力仍有待提升。
自监督学习（Self-Supervised Learning, SSL）作为解决上述挑战的有力工具，通过从无标注数据中学习丰富表征展现出巨大潜力。自监督学习的核心思想是设计巧妙的预训练任务，利用数据本身的内在结构和规律作为监督信号，从而避免对人工标注的依赖。在自然语言处理领域，双向编码器表征转换器（Bidirectional Encoder Representations from Transformers, BERT）\cite{devlin_bert_2018}通过掩码语言模型（Masked Language Model, MLM）实现了上下文感知的语言理解，能够从大规模无标注文本中学习到丰富的语义表征。生成式预训练转换器（Generative Pre-trained Transformer, GPT-3）\cite{floridi_gpt-3_2020}则通过自回归语言建模展现了强大的生成能力，在文本生成、问答和代码编写等任务中取得了令人瞩目的成果。在计算机视觉领域，掩码自编码器（Masked AutoEncoder, MAE）\cite{he2022masked}通过随机掩码图像块并重建缺失部分的方式，从大规模无标注图像中学习到了丰富的视觉表征，模型能够学习到对象的结构、纹理和语义信息。对比学习方法如简单对比学习框架SimCLR\cite{chen2020simple}和动量对比学习MoCo\cite{he2020momentum}通过最大化相似样本间的一致性、最小化不同样本间的相似性来学习判别性表征。这些成功案例充分证明了自监督学习在利用无标注数据学习丰富表征方面的强大能力，为解决标注数据稀缺问题提供了有效途径。


然而，自监督学习在EEG分析中的应用仍处于起步阶段。现有方法主要采用简单的对比学习策略\cite{krishnan_self-supervised_2022, misra_self-supervised_2020, rafiei2022self}，往往将EEG数据简单地表示为二维矩阵，忽视了电极间复杂的空间拓扑关系和神经解剖学约束。EEG信号的空间分布遵循特定的神经解剖学结构，不同脑区的功能专门化体现了大脑信息处理的层次化组织特性，而简单的矩阵表示无法有效捕获这些重要的空间语义信息。此外，现有方法通常采用简单的数据增强策略，如高斯噪声、时间偏移或频域变换\cite{shen2022contrastive}，缺乏对EEG信号特定特征的深入考虑。更重要的是，现有方法缺乏针对EEG信号噪声和伪迹问题的专门设计，容易在学习过程中将噪声模式误认为有意义的特征。


为解决上述技术挑战，本章提出了基于去噪掩码自编码器的EEG信号预训练方法DMAE-EEG（\textbf{D}enoising \textbf{M}asked \textbf{A}uto\textbf{E}ncoder for \textbf{EEG}）。该框架将神经科学先验知识与深度学习技术有机整合，具有三个核心创新点。首先，设计了脑区拓扑异质划分（Brain Region Topological Heterogeneity, BRTH）方法，基于神经解剖学先验将不规则的电极分布投影为稳定的“脑区块-补丁”图式化表征。该方法根据大脑皮层的功能分区，将电极按照脑区归属进行分组，并在每个脑区内划分为固定大小的时空补丁，实现跨设备、跨数据集的一致性表示，显著提升跨数据集对齐能力。其次，构建了去噪伪标签生成器（Denoised Pseudo-Label Generator, DPLG），通过重建型去噪任务生成稳健的教师信号。该生成器对原始EEG信号进行伪迹去除处理，生成相对干净的伪标签作为监督信号，有效缓解噪声和伪迹对预训练目标的干扰，引导模型学习具有任务无关性和可迁移性的表征。
最后，采用非对称掩码自编码器架构，利用Transformer的自注意力机制捕获长程时空依赖关系。编码器采用较深结构学习丰富表征，解码器保持轻量化设计聚焦于重建和去噪任务，平衡了表征学习的深度和计算效率。

本章将详细阐述DMAE-EEG框架的技术细节，通过在14个公开数据集上的大规模实验，验证其在信号质量增强和运动意图识别等任务中的有效性，展示其在跨会话、跨受试者和跨任务泛化方面的卓越能力。

\section{通用表征学习方法}
\subsection{DMAE-EEG预训练框架}
本文提出了一种基于去噪掩码自编码器的EEG信号预训练框架DMAE-EEG，用于构建鲁棒的EEG信号分析系统。如图~\ref{fig:C5-2}所示，DMAE-EEG通过学习可泛化的EEG表征来显著提升跨任务性能。该框架采用两阶段训练策略：第一阶段为自监督预训练阶段，开发了基于脑区拓扑异质性的分块输入策略来编码EEG信号的局部时空动态特征，设计了去噪重建任务，通过从噪声和部分观察数据中重建清洁信号来学习EEG信号的内在复杂性；第二阶段为下游任务适配阶段，将预训练编码器适配到生成式和判别式任务中，对于生成式任务，模型修复损坏的EEG信号以提高信号质量，对于判别式任务，专注于基于真实和想象运动的EEG信号进行运动识别。

\begin{figure}[htbp]
\centering
\includegraphics[width=\textwidth]{C5-2DMAE-CN.pdf}
\caption{DMAE-EEG预训练框架整体架构}
\caption*{\centering (a) 基于时空异质性的脑电分块策略；(b) 自监督预训练过程；(c) 去噪伪标签生成器；(d) 下游任务分类，包括判别式和生成式两类任务；(e) 下游任务微调，将预训练编码器适配到具体应用场景。}
\label{fig:C5-2}
\end{figure}


\subsection{脑区拓扑异质划分方法}
将基于视觉的分块技术直接应用于EEG信号面临根本性挑战，这主要源于EEG数据的内在拓扑结构特性。EEG信号通过放置在特定头皮位置的电极捕获神经活动，每个电极对应于不同的大脑功能区域。传统的图像分块方法忽略了电极配置固有的神经解剖学拓扑关系，若将来自空间上相距较远的大脑区域的信号组合在单个分块中，将违背分块应编码局部相关信息的基本前提。直接采用基于图像的分块技术未能考虑EEG信号的空间复杂性和功能连接性，存在混合不相关神经信号的风险\cite{chen2021pre, han2022survey, dosovitskiy2020image}。

为了解决上述局限性，本研究提出了一种基于脑区拓扑异质性的分块器BRTH。BRTH分块器首先基于神经解剖学先验知识，按照皮层功能区域将电极划分为$M$个局部超级节点。如图~\ref{fig:C5-2}(a)所示，根据国际10-20电极系统的标准布局，本研究定义了十个功能性超级节点，包括前额叶、运动皮层、顶叶、枕叶等关键脑区。设EEG信号段$\mathbf{x} \in \mathbb{R}^{C\times T}$具有$C$个通道和$T$个时间点，电极根据潜在的皮层区域被划分为$M$个超级节点：
\begin{equation}
\mathbf{C} = \{{C}_1,{C}_2,...,{C}_M\},\quad \sum_{m=1}^{M}|{C}_m|=C,
\end{equation}
式中，${C}_m$表示第$m$个超级节点内的电极集合，$m \in [1, M]$表示超级节点索引，$|{C}_m|$表示第$m$个超级节点包含的电极数量。为了捕获EEG信号的时间动态特性，将时间维度划分为$N$个等长度段，满足条件：
\begin{equation}
\mathbf{T} = \{T_1, T_2, ..., T_N\}, \quad \sum_{n=1}^{N}|{T}_n|=T,
\end{equation}
式中，$T_n$表示第$n$个时间段，$|{T}_n|$对应第$n$个时间段的长度，$n \in [1, N]$表示时间段索引。通过时空双重划分，每个时空分块$P_{m,n}$包含来自第$m$个超级节点在第$n$个时间段内的信号子集：
\begin{equation}
P_{m,n} = \mathbf{x}_{{C}_m,{T}_n} \in \mathbb{R}^{|{C}_m|\times |{T}_n|}, \quad m \in [1, M], n \in [1, N],
\end{equation}
式中，$P_{m,n}$表示包含第$m$个超级节点在第$n$个时间段的分块，$\mathbf{x}_{{C}_m,{T}_n}$表示从原始信号$\mathbf{x}$中提取对应区域和时间的子矩阵。这种分块策略能够有效分离来自不同大脑区域和时间段的信号，有助于捕获语义丰富且具有神经生理学意义的分块表征。



如算法~\ref{alg:graph_pooling}所示，图感知池化层为每个分块$P_{m,n}$计算$D$维嵌入$\mathbf{z}_{m,n} \in \mathbb{R}^{D}$，通过加权聚合超级节点内各通道的信息。该池化操作在保持关键神经信号特征的同时，将高维数据压缩到更紧凑的表征空间中。最终，将$M\times N$个分块级嵌入连接形成EEG令牌序列$\mathbf{z} \in \mathbb{R}^{(M\times N) \times D}$，这些令牌编码了EEG段的时空信息，为后续的表征学习奠定基础。

\begin{algorithm}[htbp]
\caption{基于图的大脑区域池化}
\label{alg:graph_pooling}

\KwIn{脑电片段 $\mathbf{x} \in \mathbb{R}^{C \times T}$，超级节点划分 $\mathbf{C} = \{C_1, C_2, ..., C_M\}$，时间划分 $\mathbf{T} = \{T_1, T_2, ..., T_N\}$，邻接矩阵 $\mathbf{A} \in \mathbb{R}^{C \times C}$ (可选)}
\KwOut{所有分块 $P_{m,n}$ 的分块嵌入 $\mathbf{z}_{m,n} \in \mathbb{R}^{D}$}
\SetKwComment{Comment}{$\triangleright$\ }{}

\If{$\mathbf{A}$ 为 \textbf{None}}{
    $\mathbf{A} \gets \mathbf{1}_{C \times C}$ \Comment{默认全连接}
}

\For{$m = 1$ \KwTo $M$}{
    $\mathbf{A}_m \gets \mathbf{A}[C_m][:, C_m]$ \Comment{提取区域连接性}
    
    $\mathbf{w} \gets \sum_{j} \mathbf{A}_m[:, j]$ \Comment{计算节点重要性权重}
    $\mathbf{w} \gets \mathbf{w} / \sum \mathbf{w}$ \Comment{权重归一化}
    
    $\mathbf{x}'[m, :] \gets \sum_{i \in C_m} \mathbf{w}_i \cdot \mathbf{x}[i, :]$ \Comment{加权通道池化}
}

\For{$m = 1$ \KwTo $M$}{
    \For{$n = 1$ \KwTo $N$}{
        $P_{m,n} \gets \mathbf{x}'[m, T_n]$ \Comment{提取时空分块}
        
        $\mathbf{z}_{m,n} \gets \text{Embedding}(P_{m,n})$ \Comment{线性投影到嵌入空间}
        
        $\mathbf{z}_{m,n} \gets \text{LayerNorm}(\mathbf{z}_{m,n})$ \Comment{层归一化}
    }
}

\Return $\{\mathbf{z}_{m,n} | m \in [1,M], n \in [1,N]\}$
\end{algorithm}



\subsection{去噪伪标签生成器}
在DMAE-EEG框架中，DPLG负责为无标签EEG信号生成高质量的伪标签，这些伪标签在预训练阶段作为自编码器的重建目标。DPLG的核心功能是通过标准化预处理流程为模型提供清洁的监督信号，同时保持数据的多样性以增强模型的噪声鲁棒性。

对于原始EEG试验记录$\mathbf{X}_{\text{trial}} \in \mathbb{R}^{C \times T_{\text{trial}}}$，式中$C$表示电极通道数，$T_{\text{trial}}$表示试验记录的时间长度，首先将其按固定窗口长度划分为$K$个EEG段。
为增强模型的噪声鲁棒性，DPLG对清洁和噪声EEG段采用差异化的预处理策略。该双重处理策略旨在确保模型在训练过程中同时接触清洁和含噪声的数据样本，从而提高对真实世界EEG记录变异性的适应能力。
对于清洁EEG信号处理流程，依次应用重采样、滤波、平均参考、基于ICA的伪影去除、信号分段和BRTH分块等完整处理步骤。重采样操作确保所有数据具有统一的采样率，带通滤波保留感兴趣的神经振荡频率范围（通常为0.5-50 Hz），平均参考去除共同模式噪声，信号分段将连续记录划分为固定长度的分析窗口。
对于噪声EEG信号处理流程，有意省略滤波、重参考和伪影去除等关键步骤，以保留真实的干扰模式和噪声特征。这种设计确保模型在自监督学习过程中能够学习到在噪声环境下的鲁棒表征能力，提高对真实应用场景的泛化性能。


对于清洁EEG信号处理流程中，为确保跨数据集的一致性和可靠性，DPLG采用标准化预处理流水线，对基于功率谱密度（Power Spectral Density, PSD）分析的自适应滤波处理和基于独立成分分析（Independent Component Analysis, ICA）的伪影去除等关键步骤进行如下分析。
为减轻电磁干扰和高频噪声的影响，首先进行PSD分析以检查EEG信号的频谱特征分布。通过分析不同频率成分的功率密度，能够有效识别系统性噪声源和异常频率成分。例如，对PhysioNet运动想象数据集的分析揭示了明显的60 Hz工频噪声峰值，如图~\ref{fig:C5-3}所示。基于频谱分析结果，应用自适应带通滤波器去除不必要的高频成分和低频漂移，同时最大限度保留神经生理学相关信息。
\begin{figure}[htbp]
\centering
\includegraphics[width=\textwidth]{C5-3freq.pdf}
\caption{电极分布和功率谱密度曲线}
\label{fig:C5-3}
\end{figure}
为进一步提升信号质量，采用基于ICA的自动伪影去除方法，利用MNE-Python工具包中的自动伪影识别算法。ICA方法基于信号统计独立性假设，能够有效地从混合的多通道EEG信号中分离眼动、肌肉收缩和运动伪影等非神经活动成分。该算法通过寻找最大化统计独立性的线性变换矩阵，将观测信号分解为相互独立的成分。伪影成分的识别基于其时域和频域特征，包括功率谱形状、时间序列的峰度和偏度等统计特征。如图~\ref{fig:C5-4}所示，ICA能够在有效去除非神经成分的同时保留重要的EEG特征。


\begin{figure}[htbp]
\centering
\includegraphics[width=0.9\textwidth]{C5-4ica-CN.pdf}
\caption{基于独立成分分析的伪影识别与剔除示意}
\caption*{\centering (a) ICA000–ICA013的空间拓扑图；(b)各成分的伪影评分，红色柱为阈值以下、被标记用于移除的成分。}
\label{fig:C5-4}
\end{figure}


经过预处理的EEG段通过BRTH分块器处理以获得嵌入表示。设第$k$个EEG段的嵌入令牌为$\mathbf{z}^k \in \mathbb{R}^{(M \times N) \times D}$，式中$M$表示空间超级节点数量，$N$表示时间分割数量，$D$表示每个分块的特征维度。单个试验记录的所有嵌入表示集合为：
\begin{equation}
\mathbf{Z}_{\text{trial}} = \{\mathbf{z}^k | \mathbf{z}^k \in \mathbb{R}^{(M \times N) \times D}, k=1, \ldots, K\},
\end{equation}
式中，$K$表示从单个试验记录中划分的EEG段总数，$k$表示EEG段索引。
DPLG生成相应的清洁伪标签集合，第$k$个EEG段的清洁标签表示为$\mathbf{y}^k \in \mathbb{R}^{(M \times N) \times D}$，对应通过完整预处理流程和BRTH分块器获得的分块级表征。清洁EEG段的最终表示为：
\begin{equation}
\mathbf{X}_{\text{clean}} = \{\mathbf{y}^k | \mathbf{y}^k \in \mathbb{R}^{(M \times N) \times D}, k=1, \ldots, K\},
\end{equation}
式中，$\mathbf{y}^k$表示第$k$个清洁EEG段的嵌入令牌。
在模型训练过程中，自编码器的输入是经过掩码操作的噪声EEG段$\mathbf{z}^k$，而相应的清洁EEG段$\mathbf{y}^k$作为重建目标提供监督信号。这种学习设计使得模型能够在噪声和部分观察条件下学习重建清洁的EEG信号表征。


\subsection{非对称编码--解码器架构}
本研究构建了一个基于Transformer的非对称编码--解码器架构\cite{he2022masked}。编码器专注于从部分可见的输入中提取丰富的时空特征信息，解码器利用这些特征在预训练期间完成信号补全和去噪重建任务。

\textbf{编码器设计。}编码器处理来自BRTH分块器的输入$\mathbf{z}^k \in \mathbb{R}^{(M\times N)\times D}$，包含$M\times N$个时空分块，每个分块具有$D$维特征表示。为模拟EEG信号采集过程中的部分观察场景，处理过程始于随机掩码操作。定义掩码向量$\mathbf{m} \in \{0,1\}^{M\times N}$，其中$m_i$表示第$i$个分块的可见性：
\begin{equation}
\begin{aligned}
&\mathbf{\hat{z}}^k_{i,j} = {m_i} \mathbf{z}^k_{i,j}, \quad  i=1,\ldots,{M\times N}, \quad j=1,\ldots,D,\\
&\text{其中 } m_i =
    \begin{cases}
    0, & \text{如果第 } i \text{ 个分块被掩码},\\
    1, & \text{如果第 } i \text{ 个分块可见}.
    \end{cases}
\end{aligned}
\end{equation}
式中，$\mathbf{\hat{z}}^k_{ij}$表示掩码后第$i$个分块的第$j$个特征分量，$i$表示分块索引，$j$表示特征维度索引。
掩码操作后，结果矩阵$\mathbf{\hat{z}}^k$包含两类分块表示：
\begin{equation}
\mathbf{\hat{z}}^k = \mathbf{\hat{z}}^k_{\text{visible}} \cup \mathbf{\hat{z}}^k_{\text{masked}},
\end{equation}
式中，$\mathbf{\hat{z}}^k_{\text{visible}}$表示当$m_i = 1$时的可见分块集合，$\mathbf{\hat{z}}^k_{\text{masked}}$表示当$m_i = 0$时的掩码分块集合。为提高训练效率，仅将$\mathbf{\hat{z}}^k_{\text{visible}}$输入到编码器中进行处理。
编码器由$L$个Transformer层组成，每层包含多头自注意力机制和前馈神经网络。对于第$l$层，多头自注意力首先计算查询、键和值投影：
\begin{equation}
\mathbf{Q}_l = \mathbf{Z}_l\mathbf{W}_l^Q, \quad
\mathbf{K}_l = \mathbf{Z}_l\mathbf{W}_l^K, \quad
\mathbf{V}_l = \mathbf{Z}_l\mathbf{W}_l^V,
\end{equation}
式中，$\mathbf{W}_l^Q, \mathbf{W}_l^K, \mathbf{W}_l^V$分别是查询、键和值投影的可学习权重矩阵，$\mathbf{Z}_l$为第$l$层的输入特征，$l \in [1, L]$表示层索引。
多头注意力机制通过并行计算多个注意力头来捕获EEG信号中的复杂时空依赖关系：
\begin{equation}
\text{MHAtt}(\mathbf{Z}_l) = \text{Concat}(\text{head}_1, ..., \text{head}_h)\mathbf{W}^O,
\end{equation}
式中，$\text{MHAtt}(\cdot)$表示多头注意力函数，$\mathbf{Z}_l$表示第$l$层的输入特征，$\text{head}_i$表示第$i$个注意力头的输出，$h$表示注意力头的数量，$\mathbf{W}^O$是输出投影矩阵。单个注意力头的计算定义为：
\begin{equation}
\text{Attention}(\mathbf{Q}_l^i, \mathbf{K}_l^i, \mathbf{V}_l^i) = \text{softmax}\left(\frac{\mathbf{Q}_l^i (\mathbf{K}_l^i)^\top}{\sqrt{d_k}}\right) \mathbf{V}_l^i,
\end{equation}
式中，$\text{Attention}(\cdot, \cdot, \cdot)$表示单头注意力计算函数；$\mathbf{Q}_l^i, \mathbf{K}_l^i, \mathbf{V}_l^i$分别表示第$l$层第$i$个注意力头的查询、键和值矩阵，$d_k$表示键向量的维度，$\text{softmax}(\cdot)$表示归一化函数。
注意力计算后，特征通过前馈神经网络进行非线性变换：
\begin{equation}
\text{FFN}(\text{MHAtt}(\mathbf{Z}_l)) = \sigma\left(\text{MHAtt}(\mathbf{Z}_l)\mathbf{W}_1 + \mathbf{b}_1\right)\mathbf{W}_2 + \mathbf{b}_2,
\end{equation}
式中，$\text{FFN}(\cdot)$表示前馈神经网络函数，$\sigma(\cdot)$表示激活函数（通常采用GELU）；$\mathbf{W}_1$和$\mathbf{W}_2$是前馈网络的权重矩阵，$\mathbf{b}_1$和$\mathbf{b}_2$是相应的偏置向量。
每层的最终输出通过残差连接和层归一化进行处理：
\begin{equation}
\mathbf{Y}_l = \text{LayerNorm}\left(\text{FFN}(\text{MHAtt}(\mathbf{Z}_l)) + \mathbf{Z}_l\right),
\end{equation}
式中，$\text{LayerNorm}(\cdot)$表示层归一化函数；$\mathbf{Y}_l$表示第$l$层的输出特征。编码器最终层的输出为编码特征表示$\mathbf{E}_k \in \mathbb{R}^{|\text{visible}| \times D_e}$，其中$D_e$是编码器的特征维度，$|\text{visible}|$表示可见分块的数量。该表示编码了从可见分块中提取的时空模式信息。


\textbf{解码器设计。}解码器接收编码器输出的特征表示$\mathbf{E}_k \in \mathbb{R}^{|\text{visible}| \times D_e}$。为保证维度兼容性，首先通过线性投影将编码特征对齐到解码器维度：
\begin{equation}
\mathbf{E}_k' = \mathbf{E}_k\mathbf{W}_{\text{align}} + b_{\text{align}},
\end{equation}
式中，$\mathbf{W}_{\text{align}} \in \mathbb{R}^{D_e \times D_d}$为投影权重矩阵，$D_d$是解码器的特征维度，$b_{\text{align}}$为偏置项。
为重建完整的EEG信号，在掩码位置插入可学习的掩码令牌$\mathbf{R}_{\text{masked}} \in \mathbb{R}^{|\text{masked}| \times D_d}$，形成完整的解码器输入：
\begin{equation}
\mathbf{Z}_k = \mathbf{E}_k' \cup \mathbf{R}_{\text{masked}} \in \mathbb{R}^{(M\times N) \times D_d},
\end{equation}
式中，$\mathbf{Z}_k$为解码器的完整输入序列，$\cup$表示按原始位置顺序重新组合可见特征和掩码令牌的操作，$|\text{masked}|$表示被掩码分块的数量。
解码器采用与编码器相似的Transformer架构，通过$K$个解码器层逐步细化掩码令牌的表示。每层解码器都能访问完整的令牌序列，使得掩码位置的预测能够充分利用可见分块的上下文信息。解码器的最终输出包含对所有位置的预测表示$\mathbf{\hat{y}}^k \in \mathbb{R}^{(M\times N) \times D}$，完成了对缺失EEG信号的重建任务。


\subsection{多组件复合损失函数}
本文提出了一种多组件复合损失函数，旨在通过协同优化重建精度与去噪能力，实现对缺失EEG信号的高质量填补。该损失函数的核心设计理念在于将信号重建与噪声抑制统一于同一优化框架，从而使模型在学习EEG信号内在结构的同时，具备对复杂噪声环境的适应能力。
复合损失函数的数学表达式为：
\begin{equation}
\mathcal{L}_{\text{dmae}} = \underbrace{||\mathbf{z}^k_{masked} - \mathbf{\hat{y}}^k_{masked}||_2}_{\text{重建损失}} + \underbrace{||\mathbf{y}^k - \mathbf{\hat{y}}^k||_2}_{\text{去噪损失}},
\end{equation}
式中，$\mathbf{\hat{y}}^k$表示自编码器对第$k$个EEG段的完整重建预测，$\mathbf{\hat{y}}^k_{masked}$表示该预测在掩码位置的取值，$\mathbf{z}^k_{masked}$表示第$k$个EEG段中掩码位置对应的原始观测信号，$\mathbf{y}^k$表示去噪伪标签生成器为第$k$个EEG段生成的清洁参考信号。

重建损失项的设计目标在于驱动模型学习EEG信号的时空依赖结构。该损失项通过约束模型在掩码位置的预测与真实观测的一致性，迫使编码器从可见信号中提取关键的神经动力学特征。这一过程本质上是一个结构化推理任务，要求模型理解EEG信号在不同导联间的空间耦合关系以及神经振荡的时域连续性。通过最小化重建损失，模型能够捕获EEG信号中蕴含的神经网络激活模式，为准确填补缺失信息奠定基础。去噪损失项的引入基于临床EEG信号固有的复杂噪声特性。实际采集的EEG数据不可避免地受到多种干扰源的污染，包括眼电伪迹、肌电噪声、工频干扰以及电极阻抗变化等。去噪损失通过将模型的整体输出与经过预处理的清洁参考信号对齐，促使模型学习区分有效神经信号与噪声成分的能力。这种设计不仅提升了模型对噪声的鲁棒性，更重要的是使模型能够在噪声环境下保持对神经活动本质特征的敏感性。


复合损失函数的协同作用机制体现在两个损失项的互补性优化过程中。重建损失确保模型维持对原始信号结构的忠实重现能力，防止过度平滑导致的信息丢失；去噪损失则引导模型学习更加抽象和鲁棒的特征表示，增强其对噪声扰动的抗干扰能力。这种双重约束机制使得训练后的编码器既能准确捕获EEG信号的核心神经动力学信息，又能有效抑制各类噪声干扰，从而为下游的EEG信号分析任务提供高质量的特征表示。

\subsection{下游任务微调}
本研究验证了预训练编码器在生成式和判别式任务中的下游应用效果。生成式任务涉及EEG信号质量增强，这可以被视为预训练过程的自然延伸。该任务的核心在于利用预训练阶段学习到的信号重建能力，对受损或低质量的EEG信号进行修复和增强。通过这种方式，模型能够将预训练阶段获得的对EEG信号内在结构和时空依赖关系的理解，直接应用于实际的信号处理任务中。

针对信号质量增强任务，本研究设计了大脑区域级别和电极级别两种不同粒度的信号缺失场景。图~\ref{fig:C5-5}展示了在25\%、50\%和75\%三种缺失率下的实验设置示意图，其中灰色区域表示被遮挡或缺失的电极位置。

\begin{figure}[htbp]
\centering
\includegraphics[width=0.95\textwidth]{C5-5gen.pdf}
\caption{信号质量增强任务中不同缺失模式示意图}
\caption*{\centering (a) 大脑区域级别缺失模式；(b) 电极级别缺失模式。}
\label{fig:C5-5}
\end{figure}

基于区域的重建任务与预训练阶段的自监督学习任务在机制上高度一致，均采用掩码策略隐藏特定脑区信号并训练模型进行重建。这种设计使得预训练知识能够无缝迁移至下游任务。相较而言，基于电极的重建任务更具实际应用价值，能够有效解决单电极故障、接触不良等实际采集场景中的信号质量问题。在微调策略上，两种重建任务均保持编码器架构不变，仅对解码器的重建模块进行参数优化。这种设计既保留了预训练阶段获得的泛化表示能力，又实现了针对特定应用场景的性能优化。


对于判别式任务验证，本研究选择运动解码这一BCI领域的经典范式。运动解码任务需要从EEG信号中准确识别和分类不同的运动意图状态，是BCI系统实现精确控制的核心技术。为充分利用预训练编码器的特征提取能力，采用编码器--分类头的架构设计，通过softmax层实现最终分类：
\begin{equation}
\hat{y}_{\text{motor}} = \text{softmax}(W_m\text{Encoder}(X_{\text{motor}}) + b_m),
\end{equation}
式中，$\hat{y}_{\text{motor}}$表示运动解码的预测输出，$X_{\text{motor}}$表示输入的运动相关EEG信号，$W_m$和$b_m$分别表示分类头的权重矩阵和偏置参数。
这种架构设计的核心思想是利用预训练编码器强大的特征提取能力，将原始EEG信号转换为高维语义特征，然后通过相对简单的分类头完成最终的分类任务。
设运动解码数据集的真实标签为$y_{\text{motor}}$，则模型优化目标可表述为：
\begin{equation}
\theta^* = \arg\min_\theta \sum_{i=1}^N L\left(y_{motor}^i, \mathrm{Classifier}\left(X_{motor}^i; \theta\right)\right),
\end{equation}
式中，$\theta^*$表示优化后的模型参数，$N$表示训练样本总数，$L$表示交叉熵损失函数，$y_{motor}^i$表示第$i$个样本的真实标签，$\mathrm{Classifier}\left(X_{motor}^i; \theta\right)$表示模型对第$i$个样本的分类预测。
在微调过程中，通常采用较低的学习率来更新预训练编码器的参数，以避免破坏预训练阶段学习到的有用表示。同时，分类头部分可以使用相对较高的学习率进行训练，因为这部分参数是随机初始化的，需要从头开始学习。这种差异化的学习率设置策略能够在保持预训练知识的同时，快速适应新的分类任务。


\section{实验设置与性能评估}
\subsection{数据集与基线方法设置}
本研究构建了一个大规模多范式EEG数据预训练框架，该框架整合了14个公开EEG数据集，系统性地覆盖了BCI领域的主要神经活动范式。数据集的选择遵循了代表性、多样性和互补性的原则，确保预训练模型能够学习到跨范式的通用EEG特征表示。
预训练数据涵盖三大核心BCI范式：运动想象范式作为最广泛应用的BCI范式，包含了BNCI2015001\cite{faller2012autocalibration}、BNCI2015004\cite{scherer2015individually}、Cho2017\cite{cho2017eeg}、Lee2019MI\cite{lee2019eeg}、MunichMI\cite{grosse2009beamforming}、Ofner2017\cite{ofner2017upper}、Schirrmeister2017\cite{schirrmeister2017deep}、Shin2017A\cite{shin2016open}、Shin2017B\cite{shin2016open}和Weibo2014\cite{yi2014evaluation}等十个数据集；事件相关电位范式通过Lee2019ERP\cite{lee2019eeg}数据集提供了外源性认知处理的神经标记；稳态视觉诱发电位范式则由Wang2016\cite{wang2016benchmark}数据集贡献了频率标记的视觉刺激响应模式。此外，PhysionetMI\cite{goldberger_physiobank_2000, schalk2004bci2000}和MultiM11\cite{jeong_multimodal_2020}数据集分别提供了大规模受试者群体的运动执行与想象数据，其中前者选取前80名受试者，后者选取前15名受试者参与预训练。表~\ref{tab:5-1}系统展示了预训练数据集的关键技术参数，这些参数的多样性构成了模型泛化能力的重要基础。受试者规模从10人扩展至109人，涵盖了小样本精细化研究到大规模群体分析的不同研究设计；通道配置从13个电极的简化布局到128个电极的高密度阵列，体现了从便携式设备到科研级设备的技术跨度；采样率从160Hz到2500Hz的宽频范围，确保了模型对不同时间分辨率数据的适应性；会话数量的变化则反映了从单次记录到纵向追踪的不同实验设计需求。

\begin{table}[htbp]
\centering
\caption{预训练数据集参数规格概览}
\label{tab:5-1}
\begin{tabular}{lcccc}
\toprule
数据集名称 & 受试者数量 & 通道数 & 采样率 (Hz) & 会话数 \\
\midrule
BNCI2015001\cite{faller2012autocalibration} & 13 & 13 & 512 & 2 \\
BNCI2015004\cite{scherer2015individually} & 10 & 30 & 256 & 2 \\
Cho2017\cite{cho2017eeg} & 53 & 64 & 512 & 1 \\
Lee2019MI\cite{lee2019eeg} & 55 & 62 & 1,000 & 2 \\
MunichMI\cite{grosse2009beamforming} & 10 & 128 & 500 & 1 \\
Ofner2017\cite{ofner2017upper} & 15 & 61 & 512 & 1 \\
Schirrmeister2017\cite{schirrmeister2017deep} & 14 & 128 & 500 & 1 \\
Shin2017A\cite{shin2016open} & 29 & 30 & 200 & 3 \\
Shin2017B\cite{shin2016open} & 29 & 30 & 200 & 3 \\
Weibo2014\cite{yi2014evaluation} & 10 & 60 & 200 & 1 \\
Lee2019ERP\cite{lee2019eeg} & 54 & 62 & 1,000 & 2 \\
Wang2016\cite{wang2016benchmark} & 32 & 62 & 250 & 1 \\
PhysionetMI$^{\dagger}$\cite{goldberger_physiobank_2000, schalk2004bci2000} & 109$^{\dagger}$ & 64 & 160 & 1 \\
MultiM11$^{\ddagger}$\cite{jeong_multimodal_2020} & 25$^{\ddagger}$ & 71 & 2,500 & 2 \\
\bottomrule
\end{tabular}
\begin{tablenotes}
\footnotesize
\item[]\centering 注：$\dagger$ 前80名受试者用于预训练, $\ddagger$ 前15名受试者用于预训练
\end{tablenotes}
\end{table}

数据预处理采用了差异化处理策略，针对清洁信号和噪声信号设计了不同的处理流程。对于清洁信号片段，实施了标准化的信号处理管道：首先通过重采样将所有数据统一至250Hz采样率，确保跨数据集的时间一致性；随后应用0.5-45Hz带通滤波器，有效抑制低频生理漂移和高频环境噪声；采用共平均参考重参考策略，消除全局共模干扰；通过独立成分分析技术自动识别并去除眼动、肌电等生理伪迹；最终将连续信号分割为2秒标准时间窗口，形成适合批处理的数据格式。
针对噪声信号片段，采用最小干预预处理策略，仅保留重采样和时间分割步骤，刻意保留了滤波前的原始伪迹特征。这种设计理念基于现实应用场景的考量，能够在真实噪声环境中进行训练，从而增强模型对实际应用中信号质量变化的鲁棒性。通过保持噪声数据的原生特性，模型能够学习识别和处理各种类型的生理与技术伪迹，这对于提升实际部署环境下的性能表现具有重要意义。
自监督预训练框架摒弃了传统监督学习对任务特定标签的依赖，仅基于原始EEG信号的内在结构进行表示学习。这种无标签学习范式能够充分利用大规模EEG数据的丰富信息，避免了标签偏置对特征表示的限制。预训练数据的多维度异质性不仅体现在神经活动范式的差异上，更涵盖了受试者的个体神经解剖差异、不同厂商设备的技术特性、实验室环境的变化以及数据采集协议的多样性。这种系统性的多样性设计使得模型能够学习到EEG信号的本质性特征，而非特定实验条件下的表面模式。



基线方法的选择遵循了历史发展脉络和技术代表性原则，构建了涵盖传统信号处理、现代深度学习和前沿预训练方法的全面对比框架。表~\ref{tab:baselines}系统梳理了不同技术路线的代表性方法及其核心特征。传统信号处理方法以共空间模式为核心，结合线性判别分析或正则化线性判别分析，代表了基于先验知识的特征工程思路；深度学习方法涵盖了卷积神经网络、注意力机制和混合架构等不同设计理念，体现了端到端学习的技术优势；预训练方法则代表了自监督学习在EEG信号处理领域的最新探索，展现了大规模无标签数据的潜在价值。


\begin{table}[htbp]
\centering
\caption{脑电信号解码基线方法设置}
\label{tab:baselines}
\begin{tabular}{lccc}
\toprule
方法 & 预训练状态 & 学习策略 & 发表年份 \\ 
\midrule
CSP+LDA\cite{wu2013common} & \texttimes & 基于CSP特征提取的线性分类 & 2013 \\
CSP+rLDA\cite{wu2014probabilistic,fu2020improvement} & \texttimes & 基于CSP特征的正则化线性分类 & 2014 \\
\midrule
FBCNet\cite{mane2020multi} & \texttimes & 多频带卷积神经网络端到端训练 & 2020 \\
EEGNet\cite{lawhernEEGnet2018} & \texttimes & 深度可分离卷积端到端训练 & 2018 \\
EEGViT\cite{arjun2021introducing} & \texttimes & 视觉变换器架构端到端训练 & 2021 \\
EEGConformer\cite{songEEG2022} & \texttimes & 卷积-注意力混合架构端到端训练 & 2022 \\
\midrule
BIOT\cite{yang2023biot} & \checkmark & 掩码自编码器预训练与下游微调 & 2023 \\
EEGPT\cite{wang2024eegpt} & \checkmark & 生成式预训练与任务适应 & 2024 \\
LaBraM\cite{jiang2024large} & \checkmark & 大规模脑信号掩码建模预训练 & 2024 \\
\bottomrule
\end{tabular}
\end{table}

信号质量增强任务采用PhysionetMI数据集受试者S81-S109的运动执行片段作为评估载体，这一选择基于该数据集在采集环境、受试者多样性和任务复杂度方面的代表性。数据分割采用3:7的微调训练与性能评估比例，确保训练样本的充分性和测试结果的可靠性。评估体系建立在多层次分析框架之上：大脑区域级别评估关注功能网络的整体恢复质量，单电极级别评估则聚焦于局部神经活动的精确重建。标准化均方误差作为点对点重建精度的定量指标，距离相关性（Distance Correlation，DC）则衡量重建信号与原始信号之间的统计依赖强度，两者结合提供了信号重建质量的全面评价。对比基线选择STC方法\cite{bahador2021reconstruction}，该方法基于时空相关性矩阵的信号重建策略代表了传统信号处理在该领域的技术水准。

运动意图识别实验构建了覆盖运动执行与想象、单肢体与多肢体、二分类与多分类的综合评估体系。PhysionetMI数据集提供了经典的手部运动执行（MM-H）和想象（MI-H）二分类任务，以及手部-足部运动执行（MM-HF）和想象（MI-HF）四分类任务，涵盖了从简单二元选择到复杂多类别决策的认知负荷梯度。MultiM11数据集进一步扩展了任务复杂度，包含手部抓握（MM-G/MI-G）的三分类任务、手臂伸展运动（MM-R/MI-R）的六分类空间定位任务，以及腕部扭转（MM-T/MI-T）的精细运动控制任务。这种多层次任务设计能够全面评估模型在不同认知复杂度和运动精度要求下的解码性能。实验采用受试者无关的跨个体泛化评估策略，训练集与测试集受试者完全分离，模拟了实际BCI应用中面临的个体差异挑战。数据分割采用7:1:2的训练、验证和测试比例，通过非重叠时间窗口确保时间独立性。性能评估指标选择平衡准确率和Cohen's Kappa系数，前者消除了类别不平衡对结果的影响，后者量化了相对于随机分类的性能提升程度。结果报告基于五次独立运行的统计分析，每次运行使用不同随机种子，通过均值和标准差展现结果的稳定性和可重复性。


\subsection{模型训练与优化参数配置}
实验环境构建在高性能分布式计算平台之上，采用Python 3.8作为开发语言，PyTorch 1.12作为深度学习框架，CUDA 11.6提供GPU加速支持。EEG信号处理依托专业库Braindecode 0.4.85和MNE 1.8.0，确保了信号处理的专业性和标准化。计算资源配置8块NVIDIA GeForce RTX 4090 GPU，采用分布式数据并行训练策略，单卡48GB显存能够支持大批次训练。

预训练阶段的超参数设计体现了自监督学习的核心思想，通过高掩码比例强化模型的表示学习能力。表~\ref{tab:5-2}展示了完整的网络架构和训练配置。视觉变换器编码器采用8层深度设计，512维嵌入空间配合16个注意力头，实现了对EEG信号时空模式的精细建模。多层感知机维度设置为1024，提供了充足的非线性变换能力。0.75的掩码比例设计基于EEG信号的冗余性特征，迫使模型学习深层的信号内在结构而非表面统计规律。训练策略采用AdamW优化器，基础学习率1e-3配合余弦退火调度，40轮预热期确保了训练初期的稳定性。512的大批次设计充分利用了分布式计算资源，加速了收敛过程。


\begin{table}[htbp]
\centering
\caption{预训练阶段网络架构与优化策略配置}
\label{tab:5-2}
\begin{tabular}{ccc}
\toprule
配置类别 & 参数名称 & 设定数值 \\
\midrule
\multirow{9}{*}{神经网络架构设计} & 时间补丁分割数量 & 40 \\
& 补丁嵌入向量维度 & 512 \\
& 编码器Transformer层数 & 8 \\
& 编码器隐藏状态维度 & 512 \\
& 多头注意力机制头数 & 16 \\
& 前馈神经网络隐藏层维度 & 1024 \\
& 随机掩码比例 & 0.75 \\
& 解码器Transformer层数 & 4 \\
& 解码器隐藏状态维度 & 512 \\
\midrule
\multirow{7}{*}{优化算法与训练策略} & 小批次样本数量 & 512 \\
& 初始基础学习率 & 1e-3 \\
& L2正则化权重衰减系数 & 5e-2 \\
& Adam优化器动量参数 & (0.9, 0.95) \\
& 余弦退火最终学习率比例 & 1e-2 \\
& 学习率线性预热轮数 & 40 \\
& 梯度下降优化算法 & AdamW \\
\bottomrule
\end{tabular}
\end{table}

微调阶段的超参数配置遵循了预训练模型保护的原则，通过相对保守的训练策略避免灾难性遗忘。表~\ref{tab:5-3}详细说明了微调过程的技术细节。批次大小降至128，减少了梯度估计的方差，提高了微调过程的稳定性。学习率设定为5e-4，相比预训练阶段显著降低，确保了对预训练权重的渐进式优化。200轮训练周期配合早停策略，在充分优化和过拟合预防之间取得平衡。不同下游任务采用差异化的参数微调策略：运动解码任务仅更新编码器和分类头参数，保持解码器权重冻结，这种策略基于运动解码任务主要依赖特征提取而非信号重建的任务特性；信号增强任务则需要端到端微调完整模型，以获得最优的重建性能。

\begin{table}[htbp]
\centering
\caption{微调阶段超参数与训练策略配置}
\label{tab:5-3}
\begin{tabular}{ccc}
\toprule
配置类别 & 参数名称 & 设定数值 \\
\midrule
\multirow{8}{*}{继承网络架构参数} & 时间补丁分割数量 & 40 \\
& 补丁嵌入向量维度 & 512 \\
& 编码器Transformer层数 & 8 \\
& 编码器隐藏状态维度 & 512 \\
& 多头注意力机制头数 & 16 \\
& 前馈神经网络隐藏层维度 & 1024 \\
& 解码器Transformer层数 & 4 \\
& 解码器隐藏状态维度 & 512 \\
\midrule
\multirow{8}{*}{微调优化策略配置} & 小批次样本数量 & 128 \\
& 最大训练迭代轮数 & 200 \\
& 微调初始学习率 & 5e-4 \\
& L2正则化权重衰减系数 & 5e-2 \\
& Adam优化器动量参数 & (0.9, 0.999) \\
& 余弦退火最终学习率比例 & 1e-2 \\
& 学习率线性预热轮数 & 5 \\
& 梯度下降优化算法 & AdamW \\
\midrule
\multirow{3}{*}{微调范围} & 运动意图解码任务 & 编码器 + 任务分类头 \\
& 补丁级信号增强任务 & 完整编码器和解码器模型 \\
& 通道级信号增强任务 & 完整模型 + 通道维度上采样层 \\
\bottomrule
\end{tabular}
\end{table}


\subsection{信号质量增强任务性能分析}

信号质量增强作为BCI系统的关键技术环节，影响着后续神经解码的准确性和系统的实用性。本节深入评估DMAE-EEG方法在不同缺失场景下的信号重建能力，并与传统时空相关性方法进行全面对比分析。图~\ref{fig:C5-6}展示了在25\%、50\%和75\%三种缺失比例条件下，两种方法在脑区级别和电极级别的重建性能对比，揭示了自监督预训练框架在信号质量增强任务中的技术优势。

\begin{figure*}[!ht]
\centering
\includegraphics[width=\textwidth]{C5-6region-CN.pdf}
\caption{不同缺失比例条件下的信号重建性能对比}
\caption*{\centering (a) 脑区级别nMSE性能；(b) 脑区级别DC性能；(c) 电极级别nMSE性能；(d) 电极级别DC性能。*** p ≤ 0.001，**** p ≤ 0.0001}
\label{fig:C5-6}
\end{figure*}

（一）脑区级信号重建性能

脑区级别的重建结果体现了两种方法在处理大尺度神经网络信号恢复任务中的根本性差异。表~\ref{tab:5-5}的详细数据揭示了多个重要的性能规律和技术特征。首先，从重建精度的角度分析，DMAE-EEG方法在所有缺失条件下均表现出显著的性能优势。在25\%缺失条件下，DMAE-EEG实现的平均nMSE为0.26±0.07，相比STC方法的0.36±0.09，实现了27.8\%的误差降低。更为重要的是，DMAE-EEG方法的标准差显著更小，表明其在不同受试者间具有更好的稳定性和一致性表现。

从个体差异的视角分析，表~\ref{tab:5-5}中不同受试者的性能变化模式揭示了重要的生理学意义。以受试者S84和S104为例，这两名受试者在STC方法下表现出相对较差的重建效果，但在DMAE-EEG方法下获得了显著改善。其中，受试者S84的nMSE指标从0.40降至0.16，S104从0.46降至0.11，改善幅度分别达到60\%和76.1\%。这种差异化改善模式表明，DMAE-EEG方法能够更好地适应个体神经解剖结构和生理特征的差异，表明大规模预训练数据中蕴含的丰富个体变异信息。

从缺失比例递增的性能退化模式分析，两种方法展现出截然不同的鲁棒性特征。STC方法随着缺失比例的增加表现出急剧的性能恶化：nMSE从25\%条件下的0.36增加到75\%条件下的0.70，增幅高达94.4\%；DC值从0.85下降至0.76，降幅达10.6\%。这种急剧退化反映了传统统计方法在处理高维稀疏数据时的根本局限性。相比之下，DMAE-EEG方法表现出卓越的抗退化能力：nMSE仅从0.26增加到0.38，增幅仅为46.2\%；DC值从0.96轻微下降至0.91，降幅仅为5.2\%。这种渐进式性能退化模式表明，掩码自编码器架构通过预训练学习到的特征表示具有更强的泛化能力和抗噪声能力。特别值得关注的是，在极端缺失条件（75\%）下，DMAE-EEG方法仍能维持接近中等缺失条件（50\%）的性能水平，这一现象在传统方法中并未观察到。深入分析受试者S89、S97和S99的数据发现，这些受试者在75\%缺失条件下的DC值甚至达到0.96，与25\%缺失条件下的最优表现相当。这种性能稳定性可能源于掩码自编码器在预训练过程中接触了大量高掩码比例的训练样本，使得模型具备了在极端稀疏条件下进行有效推理的能力。

\begin{table*}[ht]
\centering
\caption{脑区级信号重建性能的评估结果}
\label{tab:5-5}
\resizebox{1.0\textwidth}{!}{%
\begin{tabular}{ccccccccccccc}
\toprule
\multirow{3}{*}{受试者} & \multicolumn{4}{c}{掩码25\%} & \multicolumn{4}{c}{掩码50\%} & \multicolumn{4}{c}{掩码75\%} \\
\cmidrule(lr){2-5} \cmidrule(lr){6-9} \cmidrule(lr){10-13}
 & \multicolumn{2}{c}{STC} & \multicolumn{2}{c}{DMAE-EEG} & \multicolumn{2}{c}{STC} & \multicolumn{2}{c}{DMAE-EEG} & \multicolumn{2}{c}{STC} & \multicolumn{2}{c}{DMAE-EEG} \\
\cmidrule(lr){2-3} \cmidrule(lr){4-5} \cmidrule(lr){6-7} \cmidrule(lr){8-9} \cmidrule(lr){10-11} \cmidrule(lr){12-13}
 & nMSE & DC & nMSE & DC & nMSE & DC & nMSE & DC & nMSE & DC & nMSE & DC \\
\midrule
S81 & 0.44 & 0.81 & 0.24 & 0.97 & 0.46 & 0.87 & 0.33 & 0.94 & 0.67 & 0.83 & 0.33 & 0.93 \\
S82 & 0.31 & 0.85 & 0.32 & 0.97 & 0.49 & 0.83 & 0.39 & 0.94 & 0.67 & 0.80 & 0.45 & 0.91 \\
S83 & 0.36 & 0.87 & 0.31 & 0.97 & 0.43 & 0.86 & 0.31 & 0.95 & 0.68 & 0.80 & 0.41 & 0.93 \\
S84 & 0.40 & 0.76 & 0.16 & 0.94 & 0.56 & 0.75 & 0.22 & 0.87 & 0.73 & 0.71 & 0.33 & 0.85 \\
S85 & 0.53 & 0.79 & 0.37 & 0.96 & 0.70 & 0.72 & 0.34 & 0.94 & 0.71 & 0.73 & 0.48 & 0.92 \\
S86 & 0.40 & 0.88 & 0.28 & 0.96 & 0.45 & 0.85 & 0.39 & 0.92 & 0.73 & 0.78 & 0.37 & 0.91 \\
S87 & 0.26 & 0.88 & 0.35 & 0.98 & 0.46 & 0.82 & 0.41 & 0.97 & 0.70 & 0.76 & 0.45 & 0.96 \\
S89 & 0.33 & 0.90 & 0.23 & 0.98 & 0.46 & 0.88 & 0.26 & 0.97 & 0.72 & 0.79 & 0.30 & 0.96 \\
S90 & 0.35 & 0.88 & 0.29 & 0.97 & 0.57 & 0.80 & 0.41 & 0.94 & 0.71 & 0.84 & 0.43 & 0.95 \\
S91 & 0.46 & 0.88 & 0.35 & 0.96 & 0.40 & 0.90 & 0.40 & 0.94 & 0.69 & 0.88 & 0.43 & 0.88 \\
S93 & 0.45 & 0.83 & 0.31 & 0.96 & 0.56 & 0.82 & 0.37 & 0.91 & 0.78 & 0.75 & 0.34 & 0.90 \\
S94 & 0.30 & 0.88 & 0.22 & 0.98 & 0.47 & 0.81 & 0.38 & 0.95 & 0.71 & 0.75 & 0.44 & 0.93 \\
S95 & 0.17 & 0.92 & 0.15 & 0.97 & 0.46 & 0.86 & 0.26 & 0.94 & 0.53 & 0.86 & 0.35 & 0.93 \\
S96 & 0.50 & 0.83 & 0.23 & 0.96 & 0.45 & 0.87 & 0.33 & 0.92 & 0.74 & 0.73 & 0.37 & 0.90 \\
S97 & 0.37 & 0.90 & 0.26 & 0.97 & 0.42 & 0.88 & 0.31 & 0.95 & 0.59 & 0.83 & 0.33 & 0.94 \\
S98 & 0.32 & 0.86 & 0.27 & 0.97 & 0.48 & 0.82 & 0.39 & 0.94 & 0.64 & 0.76 & 0.42 & 0.94 \\
S99 & 0.27 & 0.85 & 0.34 & 0.98 & 0.51 & 0.79 & 0.40 & 0.96 & 0.66 & 0.76 & 0.40 & 0.96 \\
S101 & 0.36 & 0.81 & 0.30 & 0.97 & 0.50 & 0.79 & 0.39 & 0.93 & 0.75 & 0.77 & 0.38 & 0.93 \\
S102 & 0.43 & 0.86 & 0.31 & 0.97 & 0.59 & 0.84 & 0.40 & 0.94 & 0.70 & 0.80 & 0.43 & 0.84 \\
S103 & 0.20 & 0.88 & 0.19 & 0.96 & 0.36 & 0.84 & 0.31 & 0.93 & 0.67 & 0.56 & 0.29 & 0.93 \\
S104 & 0.46 & 0.82 & 0.11 & 0.95 & 0.54 & 0.81 & 0.21 & 0.91 & 0.79 & 0.53 & 0.26 & 0.89 \\
S105 & 0.36 & 0.88 & 0.19 & 0.93 & 0.49 & 0.81 & 0.27 & 0.92 & 0.68 & 0.82 & 0.35 & 0.91 \\
S106 & 0.37 & 0.84 & 0.27 & 0.97 & 0.56 & 0.79 & 0.42 & 0.93 & 0.74 & 0.74 & 0.46 & 0.91 \\
S107 & 0.30 & 0.89 & 0.28 & 0.97 & 0.52 & 0.77 & 0.38 & 0.94 & 0.64 & 0.83 & 0.32 & 0.93 \\
S108 & 0.27 & 0.89 & 0.32 & 0.95 & 0.57 & 0.81 & 0.38 & 0.93 & 0.71 & 0.83 & 0.43 & 0.83 \\
S109 & 0.46 & 0.75 & 0.24 & 0.94 & 0.67 & 0.66 & 0.25 & 0.86 & 0.79 & 0.47 & 0.32 & 0.84 \\
\midrule
均值 & 0.36 & 0.85 & 0.26 & 0.96 & 0.50 & 0.82 & 0.34 & 0.93 & 0.70 & 0.76 & 0.38 & 0.91 \\
标准差 & 0.09 & 0.04 & 0.07 & 0.01 & 0.08 & 0.05 & 0.06 & 0.02 & 0.06 & 0.10 & 0.06 & 0.04 \\
\bottomrule
\end{tabular}%
}
\end{table*}

（二）电极级信号重建性能

电极级别的重建任务相比脑区级别面临更为严峻的技术挑战，主要体现在空间分辨率的细致性、噪声敏感性的增强以及局部特征建模的复杂性。表~\ref{tab:5-6}的详细结果揭示了在这一更为苛刻的测试条件下，两种方法的性能特征和差异化表现模式。
从平均性能水平分析，DMAE-EEG方法在各缺失条件下均保持显著优势，但优势幅度相比脑区级别有所收窄，这符合理论预期。在25\%缺失条件下，DMAE-EEG的nMSE为0.36±0.06，相比STC的0.51±0.09，实现了29.4\%的误差降低；在75\%极端条件下，DMAE-EEG的nMSE为0.39±0.07，相比STC的0.78±0.04，误差降低幅度达到50.0\%。这种在极端条件下优势放大的现象表明，深度学习方法的技术优势在面临更大挑战时更加凸显。


\begin{table*}[htbp]
\centering
\caption{电极级信号重建性能的精细化评估结果}
\label{tab:5-6}
\resizebox{1.0\textwidth}{!}{%
\begin{tabular}{ccccccccccccc}
\toprule
\multirow{3}{*}{受试者} & \multicolumn{4}{c}{掩码25\%} & \multicolumn{4}{c}{掩码50\%} & \multicolumn{4}{c}{掩码75\%} \\
\cmidrule(lr){2-5} \cmidrule(lr){6-9} \cmidrule(lr){10-13}
 & \multicolumn{2}{c}{STC} & \multicolumn{2}{c}{DMAE-EEG} & \multicolumn{2}{c}{STC} & \multicolumn{2}{c}{DMAE-EEG} & \multicolumn{2}{c}{STC} & \multicolumn{2}{c}{DMAE-EEG} \\
\cmidrule(lr){2-3} \cmidrule(lr){4-5} \cmidrule(lr){6-7} \cmidrule(lr){8-9} \cmidrule(lr){10-11} \cmidrule(lr){12-13}
 & nMSE & DC & nMSE & DC & nMSE & DC & nMSE & DC & nMSE & DC & nMSE & DC \\
\midrule
S81 & 0.47 & 0.79 & 0.24 & 0.94 & 0.57 & 0.77 & 0.22 & 0.94 & 0.76 & 0.76 & 0.26 & 0.93 \\
S82 & 0.51 & 0.78 & 0.34 & 0.89 & 0.60 & 0.76 & 0.43 & 0.87 & 0.77 & 0.74 & 0.44 & 0.87 \\
S83 & 0.42 & 0.80 & 0.45 & 0.89 & 0.54 & 0.79 & 0.46 & 0.89 & 0.74 & 0.77 & 0.47 & 0.89 \\
S84 & 0.69 & 0.65 & 0.35 & 0.83 & 0.75 & 0.64 & 0.34 & 0.78 & 0.85 & 0.63 & 0.34 & 0.76 \\
S85 & 0.67 & 0.68 & 0.47 & 0.89 & 0.74 & 0.66 & 0.42 & 0.90 & 0.85 & 0.64 & 0.37 & 0.92 \\
S86 & 0.46 & 0.79 & 0.40 & 0.87 & 0.60 & 0.76 & 0.42 & 0.86 & 0.76 & 0.75 & 0.43 & 0.83 \\
S87 & 0.50 & 0.74 & 0.49 & 0.89 & 0.60 & 0.73 & 0.48 & 0.90 & 0.77 & 0.71 & 0.44 & 0.89 \\
S89 & 0.46 & 0.82 & 0.27 & 0.93 & 0.57 & 0.81 & 0.25 & 0.93 & 0.76 & 0.80 & 0.26 & 0.91 \\
S90 & 0.44 & 0.81 & 0.42 & 0.87 & 0.54 & 0.79 & 0.46 & 0.91 & 0.76 & 0.78 & 0.46 & 0.89 \\
S91 & 0.47 & 0.83 & 0.38 & 0.93 & 0.54 & 0.80 & 0.35 & 0.90 & 0.78 & 0.80 & 0.44 & 0.90 \\
S93 & 0.59 & 0.75 & 0.27 & 0.87 & 0.68 & 0.72 & 0.38 & 0.87 & 0.81 & 0.71 & 0.39 & 0.87 \\
S94 & 0.61 & 0.75 & 0.42 & 0.90 & 0.68 & 0.72 & 0.34 & 0.91 & 0.82 & 0.69 & 0.37 & 0.89 \\
S95 & 0.25 & 0.86 & 0.35 & 0.84 & 0.42 & 0.85 & 0.37 & 0.82 & 0.66 & 0.83 & 0.33 & 0.83 \\
S96 & 0.54 & 0.78 & 0.41 & 0.88 & 0.63 & 0.75 & 0.43 & 0.89 & 0.81 & 0.73 & 0.47 & 0.82 \\
S97 & 0.47 & 0.81 & 0.32 & 0.89 & 0.54 & 0.79 & 0.28 & 0.92 & 0.76 & 0.77 & 0.39 & 0.90 \\
S98 & 0.52 & 0.75 & 0.31 & 0.94 & 0.62 & 0.73 & 0.37 & 0.91 & 0.78 & 0.71 & 0.36 & 0.92 \\
S99 & 0.44 & 0.77 & 0.40 & 0.87 & 0.56 & 0.75 & 0.53 & 0.87 & 0.75 & 0.73 & 0.54 & 0.87 \\
S101 & 0.57 & 0.73 & 0.40 & 0.89 & 0.65 & 0.71 & 0.41 & 0.88 & 0.81 & 0.68 & 0.45 & 0.86 \\
S102 & 0.50 & 0.79 & 0.27 & 0.91 & 0.58 & 0.77 & 0.30 & 0.92 & 0.77 & 0.75 & 0.35 & 0.91 \\
S103 & 0.40 & 0.79 & 0.34 & 0.83 & 0.52 & 0.78 & 0.32 & 0.85 & 0.71 & 0.76 & 0.34 & 0.84 \\
S104 & 0.55 & 0.76 & 0.28 & 0.85 & 0.62 & 0.73 & 0.35 & 0.83 & 0.80 & 0.71 & 0.31 & 0.80 \\
S105 & 0.56 & 0.80 & 0.32 & 0.89 & 0.66 & 0.77 & 0.32 & 0.88 & 0.80 & 0.75 & 0.31 & 0.86 \\
S106 & 0.63 & 0.73 & 0.38 & 0.86 & 0.71 & 0.71 & 0.42 & 0.85 & 0.84 & 0.69 & 0.42 & 0.83 \\
S107 & 0.49 & 0.80 & 0.40 & 0.84 & 0.60 & 0.80 & 0.39 & 0.86 & 0.76 & 0.78 & 0.41 & 0.85 \\
S108 & 0.48 & 0.80 & 0.39 & 0.87 & 0.58 & 0.78 & 0.41 & 0.89 & 0.76 & 0.76 & 0.48 & 0.78 \\
S109 & 0.63 & 0.71 & 0.35 & 0.84 & 0.76 & 0.66 & 0.45 & 0.77 & 0.87 & 0.66 & 0.41 & 0.79 \\
\midrule
均值 & 0.51 & 0.77 & 0.36 & 0.88 & 0.61 & 0.75 & 0.38 & 0.88 & 0.78 & 0.73 & 0.39 & 0.86 \\
标准差 & 0.09 & 0.04 & 0.06 & 0.03 & 0.08 & 0.05 & 0.07 & 0.04 & 0.04 & 0.05 & 0.07 & 0.05 \\
\bottomrule
\end{tabular}%
}
\end{table*}

从个体差异性分析的角度，电极级重建结果展现了更为复杂的性能分布模式。例如，受试者S84、S85在STC方法下表现出较差的重建效果，在25\%信号缺失条件下，nMSE值分别达到0.69和0.67，但在DMAE-EEG方法下获得显著改善，分别降至0.35和0.47，改善幅度达到49.3\%和29.9\%。这种个体化的性能改善模式反映了深度学习方法在处理个体神经解剖差异和生理变异性方面的技术优势。受试者S99在不同缺失条件下表现出异常模式，该受试者在DMAE-EEG方法下，随着缺失比例从25\%增加到75\%，其nMSE值从0.40增加到0.54，增幅达到35\%，明显高于其他受试者。这种异常模式可能反映了特定个体的神经信号特征与预训练数据分布的匹配度差异，提示了个体化模型优化的重要性。

从方法稳定性的角度分析，DMAE-EEG在电极级重建任务中表现出更优的跨条件一致性。其标准差在各缺失条件下均维持在相对较低的水平，而STC方法的标准差相对较高，这种差异表明深度学习方法具有更强的泛化能力和更稳定的性能表现。
综合性能分析表明，DMAE-EEG方法不仅在数值指标上显著优于传统方法，更重要的是展现出在极端条件下的技术鲁棒性和实用价值。这种性能优势为BCI系统在信号质量受损情况下的可靠运行提供了重要的技术保障。


\subsection{运动意图识别任务性能评估}

运动意图识别作为BCI系统的核心功能，其性能直接决定了系统的实用性和用户体验质量。本节通过在PhysioNetMI和MultiM11两个权威基准数据集上的系统性实验，全面评估DMAE-EEG方法在不同复杂度运动解码任务中的性能表现，并通过精心设计的消融实验深入剖析各核心组件对整体性能的贡献机制。


（一）与基线方法的性能比较

表~\ref{tab:C5-7}和表~\ref{tab:C5-8}系统展示了DMAE-EEG方法在PhysioNetMI数据集上的性能表现。该数据集涵盖运动执行（MM）和运动想象（MI）两种神经活动范式，包含手部运动（H）的二分类任务和手足联合运动（HF）的四分类任务，构成了从简单到复杂的任务难度梯度。DMAE-EEG方法在所有评估任务中均实现了最优性能，其中在MI-H任务中取得85.94±0.97\%的平衡准确率和0.72±0.03的Cohen's Kappa评分，相较于最佳非预训练方法EEGConformer分别实现了3.16\%的准确率提升和9.1\%的Kappa系数改善，充分展现了自监督预训练框架的技术优势。

\begin{table*}[htbp]
\centering
\caption{PhysioNetMI数据集平衡准确率对比}
\label{tab:C5-7}
\begin{tabular}{cccccc}
\toprule
\multirow{2}{*}{方法} & \multirow{2}{*}{PT} & \multicolumn{4}{c}{平衡准确率 (\%)} \\ 
\cmidrule(lr){3-6} 
 &  & MM-H & MM-HF & MI-H & MI-HF \\ 
\midrule
CSP+LDA & \texttimes & 75.83±0.73 & 48.06±0.98 & 74.17±1.19 & 44.56±0.73 \\
CSP+rLDA & \texttimes & 75.83±1.68 & 49.31±1.53 & 74.72±1.62 & 44.39±0.81 \\
FBCNet & \texttimes & 72.50±0.60 & 65.20±0.55 & 74.00±0.65 & 68.80±0.70 \\
EEGNet & \texttimes & 73.96±0.04 & 66.93±0.57 & 75.17±0.63 & 68.75±0.71 \\
EEGViT & \texttimes & 71.53±2.62 & \underline{68.49±3.51} & 77.95±1.82 & \underline{69.01±1.56} \\
EEGConformer & \texttimes & \underline{79.44±0.88} & 68.33±0.59 & \underline{82.78±0.87} & 55.56±1.52 \\ 
\midrule
EEGPT & \checkmark & 78.89±2.81 & 55.28±1.47 & 74.44±1.67 & 54.72±3.54 \\
BIOT & \checkmark & 71.11±1.63 & 68.89±2.54 & 71.67±1.66 & 55.83±2.49 \\
LaBraM & \checkmark & 80.86±0.94 & 71.67±0.63 & 83.33±0.93 & 71.11±0.67 \\
DMAE-EEG & \checkmark & \textbf{81.08±0.26} & \textbf{72.66±0.73} & \textbf{85.94±0.97} & \textbf{72.40±0.95} \\ 
\bottomrule
\end{tabular}
\begin{tablenotes}
\footnotesize
\item[]\centering 注：非预训练方法最佳结果用\underline{下划线}标识，预训练方法最佳结果用\textbf{粗体}标识。
\end{tablenotes}
\end{table*}

\begin{table*}[htbp]
\centering
\caption{PhysioNetMI数据集Cohen's Kappa系数对比}
\label{tab:C5-8}
\begin{tabular}{cccccc}
\toprule
\multirow{2}{*}{方法} & \multirow{2}{*}{PT} & \multicolumn{4}{c}{Cohen's Kappa} \\ 
\cmidrule(lr){3-6} 
 &  & MM-H & MM-HF & MI-H & MI-HF \\ 
\midrule
CSP+LDA & \texttimes & 0.52±0.02 & 0.31±0.03 & 0.49±0.03 & 0.26±0.02 \\
CSP+rLDA & \texttimes & 0.52±0.04 & 0.32±0.04 & 0.50±0.04 & 0.26±0.04 \\
FBCNet & \texttimes & 0.46±0.02 & 0.54±0.02 & 0.49±0.02 & 0.58±0.02 \\
EEGNet & \texttimes & 0.48±0.00 & 0.56±0.02 & 0.51±0.02 & 0.58±0.02 \\
EEGViT & \texttimes & 0.44±0.06 & \underline{0.58±0.08} & 0.55±0.04 & \underline{0.59±0.05} \\
EEGConformer & \texttimes & \underline{0.59±0.03} & 0.57±0.02 & \underline{0.66±0.03} & 0.41±0.04 \\ 
\midrule
EEGPT & \checkmark & 0.58±0.06 & 0.40±0.04 & 0.49±0.04 & 0.39±0.07 \\
BIOT & \checkmark & 0.42±0.04 & 0.58±0.05 & 0.43±0.04 & 0.41±0.06 \\
LaBraM & \checkmark & 0.62±0.03 & 0.62±0.01 & 0.66±0.03 & 0.61±0.02 \\
DMAE-EEG & \checkmark & \textbf{0.62±0.01} & \textbf{0.63±0.02} & \textbf{0.72±0.03} & \textbf{0.63±0.03} \\ 
\bottomrule
\end{tabular}
\begin{tablenotes}
\footnotesize
\item[]\centering 注：非预训练方法最佳结果用\underline{下划线}标识，预训练方法最佳结果用\textbf{粗体}标识。
\end{tablenotes}
\end{table*}

从任务复杂度对性能影响的深度分析中发现，二分类任务（MM-H, MI-H）与四分类任务（MM-HF, MI-HF）之间呈现显著的性能分化特征。传统CSP方法在面临四分类挑战时表现出严重的性能坍塌，准确率普遍跌落至50\%以下，逼近随机分类的理论下限，这一现象深刻反映了传统线性方法在处理高维非线性决策边界时的根本性局限。相比之下，深度学习架构展现出更强的复杂模式建模能力，其中DMAE-EEG在最具挑战性的MI-HF任务中仍能维持72.40\%的准确率，相较于传统方法实现了62.8\%的性能跃升。
预训练方法间的横向对比分析进一步揭示了不同技术路线的内在差异。EEGPT作为生成式预训练的代表性方法，在复杂多分类任务中出现明显的性能退化，这主要源于其生成式预训练目标与判别式下游任务之间的目标函数不匹配问题。BIOT方法虽然采用了掩码自编码器架构，但受限于相对较小的预训练数据规模和单一模态设计，其性能表现仍显著逊色于DMAE-EEG。LaBraM作为当前大规模脑信号预训练的标杆性工作，在各项任务中均表现出良好的稳定性，但仍被DMAE-EEG方法全面超越，平均准确率改进幅度达到1.53\%，Cohen's Kappa系数提升4.8\%。


表~\ref{tab:C5-9}和表~\ref{tab:C5-10}进一步验证了DMAE-EEG方法在更为复杂的MultiM11数据集上的卓越泛化能力。该数据集涵盖三种不同的精细运动模式：手部抓握（G）、手臂空间到达（R）和腕部扭转（T），构成了从三分类到六分类的全方位任务复杂度测试平台。DMAE-EEG方法在全部六个评估任务中均实现最优性能，在相对简单的MM-G任务中达到89.78\%的准确率，在最具挑战性的六分类空间定位任务MM-R中仍能实现57.85\%的准确率，显著超越同类方法。

\begin{table*}[htbp]
\centering
\caption{MultiM11数据集平衡准确率对比}
\label{tab:C5-9}
\resizebox{\textwidth}{!}{%
\begin{tabular}{cccccccc}
\toprule
\multirow{2}{*}{方法} & \multirow{2}{*}{PT} & \multicolumn{6}{c}{平衡准确率 (\%)} \\ 
\cmidrule(lr){3-8} 
 &  & MM-G & MM-R & MM-T & MI-G & MI-R & MI-T \\ 
\midrule
CSP+LDA & \texttimes & 85.56±0.82 & 41.44±0.46 & 72.67±0.73 & 79.78±0.76 & 33.89±0.38 & \underline{77.33±0.83} \\
CSP+rLDA & \texttimes & \underline{86.89±0.79} & 40.22±0.43 & \underline{76.67±0.81} & 81.78±0.84 & 33.11±0.36 & 77.00±0.78 \\
FBCNet & \texttimes & 78.40±0.70 & 54.15±0.56 & 75.80±0.75 & 79.60±0.73 & 48.85±0.53 & 76.20±0.72 \\
EEGNet & \texttimes & 79.33±0.68 & 52.29±0.54 & 73.67±0.72 & 78.89±0.71 & 47.01±0.49 & 73.67±0.69 \\
EEGViT & \texttimes & 83.78±2.97 & 54.24±0.87 & 76.67±0.95 & 81.33±0.82 & 47.99±0.92 & 75.67±0.94 \\
EEGConformer & \texttimes & 85.20±0.86 & \underline{56.00±0.59} & 76.20±0.79 & \underline{84.67±0.88} & \underline{52.87±0.55} & 74.60±0.71 \\ 
\midrule
EEGPT & \checkmark & 84.67±1.83 & 56.07±0.75 & 81.20±0.91 & 82.13±1.85 & 51.27±0.57 & 79.40±0.89 \\
BIOT & \checkmark & 79.73±2.71 & 53.53±0.55 & 77.80±0.85 & 84.00±0.87 & 49.40±0.53 & 76.40±0.82 \\
LaBraM & \checkmark & 83.60±0.84 & 52.80±0.54 & 81.20±0.43 & 85.87±0.96 & 52.60±0.56 & 79.60±0.88 \\
DMAE-EEG & \checkmark & \textbf{89.78±0.45} & \textbf{57.85±0.63} & \textbf{82.33±0.97} & \textbf{87.11±0.53} & \textbf{53.33±0.58} & \textbf{80.33±0.92} \\ 
\bottomrule
\end{tabular}%
}
\begin{tablenotes}
\footnotesize
\item[]\centering 注：非预训练方法最佳结果用\underline{下划线}标识，预训练方法最佳结果用\textbf{粗体}标识。
\end{tablenotes}
\end{table*}


\begin{table*}[htbp]
\centering
\caption{MultiM11数据集Cohen's Kappa系数对比}
\label{tab:C5-10}
\begin{tabular}{cccccccc}
\toprule
\multirow{2}{*}{方法} & \multirow{2}{*}{PT} & \multicolumn{6}{c}{Cohen's Kappa} \\ 
\cmidrule(lr){3-8} 
 &  & MM-G & MM-R & MM-T & MI-G & MI-R & MI-T \\ 
\midrule
CSP+LDA & \texttimes & 0.78±0.03 & 0.30±0.02 & 0.44±0.03 & 0.69±0.03 & 0.21±0.01 & \underline{0.56±0.04} \\
CSP+rLDA & \texttimes & \underline{0.80±0.02} & 0.28±0.02 & \underline{0.52±0.03} & 0.72±0.04 & 0.20±0.01 & 0.56±0.03 \\
FBCNet & \texttimes & 0.67±0.03 & 0.42±0.02 & 0.50±0.03 & 0.69±0.03 & 0.40±0.02 & 0.54±0.03 \\
EEGNet & \texttimes & 0.69±0.03 & 0.38±0.02 & 0.47±0.03 & 0.68±0.03 & 0.36±0.02 & 0.47±0.03 \\
EEGViT & \texttimes & 0.75±0.12 & 0.44±0.05 & 0.52±0.06 & 0.72±0.04 & 0.37±0.05 & 0.51±0.06 \\
EEGConformer & \texttimes & 0.78±0.03 & \underline{0.47±0.02} & 0.52±0.03 & \underline{0.77±0.05} & \underline{0.43±0.02} & 0.49±0.04 \\ 
\midrule
EEGPT & \checkmark & 0.77±0.09 & 0.47±0.05 & 0.62±0.05 & 0.73±0.09 & 0.42±0.02 & 0.59±0.04 \\
BIOT & \checkmark & 0.70±0.11 & 0.44±0.02 & 0.56±0.04 & 0.76±0.04 & 0.39±0.02 & 0.53±0.03 \\
LaBraM & \checkmark & 0.75±0.03 & 0.43±0.02 & 0.62±0.01 & 0.79±0.06 & 0.43±0.02 & 0.59±0.05 \\
DMAE-EEG & \checkmark & \textbf{0.84±0.02} & \textbf{0.48±0.03} & \textbf{0.64±0.03} & \textbf{0.81±0.02} & \textbf{0.43±0.02} & \textbf{0.60±0.06} \\ 
\bottomrule
\end{tabular}
\begin{tablenotes}
\footnotesize
\item[]\centering 注：非预训练方法最佳结果用\underline{下划线}标识，预训练方法最佳结果用\textbf{粗体}标识。
\end{tablenotes}
\end{table*}

六分类手臂伸展任务（MM-R/MI-R）要求区分六个不同的空间方向，对模型的空间表征能力提出了极高要求。DMAE-EEG在此任务上的准确率也相对较低（MM-R: 57.85\%，MI-R: 53.33\%），但仍显著优于其他方法。这反映了当前EEG解码技术在处理精细空间信息时的技术挑战。从方法稳定性的统计学角度进行深入分析，DMAE-EEG展现出卓越的跨实验一致性表现。在MultiM11数据集的绝大多数评估任务中，DMAE-EEG的性能标准差均显著低于同类方法，如在MM-G任务中的0.45\%相较于其他方法的0.68\%-2.97\%范围，这种低方差特性表明该方法具备更强的算法鲁棒性和结果可重现性，对于实际BCI系统的可靠部署具有重要意义。


（二）核心组件贡献的消融分析

为深入解析DMAE-EEG方法各核心组件的技术贡献机制，本研究设计了全面的消融实验体系，系统性地考察三个关键设计维度：分块分辨率配置、空间嵌入策略选择以及自监督预训练任务组合。表~\ref{tab:C5-11}详细展示了消融实验的量化结果，为算法的技术优化和理论理解提供了重要的实证依据。

\begin{table}[htbp]
\centering
\caption{DMAE-EEG核心组件消融的平衡准确率}
\label{tab:C5-11}
\begin{tabular}{cccccccc}
\toprule
\multirow{2}{*}{配置类型} & \multirow{2}{*}{\makecell{分块\\数量}} & \multirow{2}{*}{\makecell{分块\\策略}} & \multirow{2}{*}{\makecell{预训练\\任务}} & \multicolumn{2}{c}{PhysioNetMI} & \multicolumn{2}{c}{MultiM11} \\ 
\cmidrule(lr){5-6} \cmidrule(lr){7-8} 
 &  &  &  & MM & MI & MM & MI \\ 
\midrule
DMAE-EEG & 40 & BRTH & MA+DE & \textbf{76.87} & \textbf{79.17} & \textbf{76.65} & \textbf{73.59} \\ 
\midrule
\multirow{3}{*}{分块分辨率变化} 
 & 20 & BRTH & MA+DE & 72.92 & 74.48 & 74.14 & 70.83 \\ 
 & 60 & BRTH & MA+DE & \underline{76.39} & \underline{78.47} & \underline{76.22} & \underline{72.22} \\ 
 & 40 & 随机 & MA+DE & 75.00 & 74.31 & 71.53 & 68.23 \\ 
\midrule
\multirow{2}{*}{预训练任务简化}  
 & 40 & BRTH & MA & 73.67 & 72.44 & 70.32 & 69.27 \\  
 & 40 & BRTH & DE & 72.24 & 70.67 & 71.53 & 67.71 \\  
\bottomrule
\end{tabular}
\begin{tablenotes}
\footnotesize
\item[]\centering 注：MA表示掩码自编码任务，DE表示去噪编码任务。最佳结果用\textbf{粗体}标示，次佳结果用\underline{下划线}标示。
\end{tablenotes}
\end{table}

分块分辨率对解码性能的影响呈现明显的非线性优化曲线特征。当分块数量从最优配置的40个显著减少至20个时，PhysioNetMI数据集上运动执行和运动想象任务的平均准确率分别下降3.95\%和4.69\%，MultiM11数据集上的性能退化幅度为2.51\%和2.76\%。这种性能衰减主要归因于信息粒度损失：过于粗糙的分块分辨率导致EEG信号中蕴含的精细时空神经活动模式无法被有效捕获和表征。相反，当分块数量增加至60个时，虽然理论上能够保留更多细节信息，但性能提升呈现边际递减效应，同时计算复杂度呈二次方增长，表明系统达到了信息冗余的临界点。

分块策略选择对算法性能的影响深刻体现了神经科学领域知识在算法设计中的关键价值。基于脑区功能解剖结构的BRTH策略相较于随机分块方法，在PhysioNetMI数据集的运动执行和想象任务上分别实现了1.87\%和4.86\%的平衡准确率提升，在MultiM11数据集上的改进效果更为显著，分别达到5.12\%和5.36\%。由此可见，遵循大脑皮层功能区划分的分块策略能够更好地保持相关神经元群体之间的协同活动模式，有利于模型学习到符合神经生理机制的有意义特征表示，而随机分块则可能破坏这种内在的神经网络组织结构。

预训练任务配置的消融分析证实了掩码自编码与去噪任务的协同效应。单独使用掩码自编码时，PhysioNetMI数据集的运动执行和想象任务上的性能下降3.20\%和6.73\%，MultiM11数据集的运动执行和想象任务上分别下降6.33\%和4.32\%。单独使用去噪任务（DE）的性能退化更为严重，在PhysioNetMI的运动执行和想象任务上分别下降4.63\%和8.50\%，在MultiM11上的运动执行和想象任务上分别下降5.12\%和5.88\%。这种互补性源于两种任务的不同学习目标，掩码自编码促进模型学习信号的全局表示，而去噪任务增强模型对局部噪声的鲁棒性。

消融实验的整体分析表明，DMAE-EEG方法的技术优势源于多个组件的精心设计和有机结合。40个分块的配置在计算效率和信息保留之间达到最佳平衡，BRTH分块策略有效利用了神经科学先验知识，掩码自编码与去噪任务的联合预训练实现了特征学习的互补增强。这些设计选择的有效性在不同数据集和任务复杂度条件下得到了一致验证，为EEG信号处理领域的方法设计提供了重要参考。

\section{结果讨论}

\subsection{脑电信号重构的鲁棒性}
DMAE-EEG模型在不同程度的数据损坏条件下均表现出卓越的信号重构鲁棒性，特别是在传统方法通常失效的高缺失比率场景中。这种鲁棒性主要源于模型采用的高掩码率自监督预训练策略，该策略使模型具备了在局部信息不足时充分利用全局上下文信息的能力。预训练过程使DMAE-EEG能够学习到EEG信号模式的可泛化表征，从而构建出可在重构过程中应用的神经动力学内部模型。
在25\%缺失比率下，模型重构精度较高，这主要归因于剩余75\%的信号信息足以支撑高精度的模式识别和预测。当缺失比率提升至50\%时，重构性能虽有所下降，但仍维持在可接受范围内，这表明模型已有效学习到EEG信号的冗余结构和关键特征分布。在极端的75\%缺失条件下，DMAE-EEG相较于传统方法仍保持显著优势，这一现象揭示了深度学习方法在处理高维时序数据不完整性问题上的独特价值。

图~\ref{fig:C5-7}所展示的电极级重构效果验证了DMAE-EEG在空间维度上的建模能力。通过对来自未参与预训练的受试者数据进行可视化分析，可以观察到模型在重构过程中不仅保持了单个电极的时域特征，更重要的是维持了电极间的空间相关性模式。这种多尺度建模能力对于保证重构信号的神经生理学意义至关重要，特别是在需要精确空间定位的临床应用中，如癫痫病灶定位和认知功能评估等场景。模型在处理不同时间窗口（0-0.5秒至1.5-2秒）时表现出的一致性，进一步证实了其对EEG信号时间动态特性的有效捕获能力。


\begin{figure*}[htbp]
\centering
\includegraphics[width=\textwidth]{C5-7vis.pdf}
\caption{脑电片段的信号重构效果对比分析}
\caption*{\centering 注：GT表示真实信号，列表示25\%、50\%、75\%缺失比率下的重构结果，行表示四个连续0.5秒时间窗口，黄色区域标示关键对比区域。}
\label{fig:C5-7}
\end{figure*}

在技术层面上，DMAE-EEG的鲁棒性表现还体现在其对不同类型噪声和伪迹的适应能力上。传统的信号处理方法往往针对特定类型的噪声设计，而DMAE-EEG通过端到端的学习框架，能够自适应地识别和补偿多种类型的信号退化。这种适应性使得该方法在实际应用中具有更强的实用性，因为真实环境中的EEG信号往往同时受到多种因素的影响。
% 从计算效率的角度来看，DMAE-EEG在推理阶段的计算复杂度相对较低，使其能够在实时或近实时的应用场景中部署。这一特性对于需要即时反馈的临床监测系统和BCI应用尤为重要。同时，模型的轻量化设计使其能够在计算资源有限的移动设备或便携式监测设备上运行，进一步扩展了其应用潜力。

\subsection{跨会话泛化能力的长期稳定性}
BCI系统在实际部署中面临的核心挑战之一是会话间信号分布的非平稳性，这种变异性源于电极阻抗漂移、头皮--电极接触状态变化以及受试者生理状态波动等多重因素。为量化评估DMAE-EEG在此类现实条件下的性能稳定性，本研究基于MultiM11数据集设计了严格的跨会话验证协议。实验采用受试者特异性的分割策略：将未参与预训练的受试者的首次会话数据作为微调基准，后续两次会话作为独立测试集，从而模拟BCI系统长期使用过程中的性能退化与恢复场景。实验设计中的自适应校准机制采用渐进式数据量配置，从10个片段（25秒）递增至50个片段（125秒），每个片段时长为2.5秒。这一设计旨在评估模型在不同校准数据预算约束下的性能表现，为实际应用中的校准策略提供定量指导。图~\ref{fig:5-8}展示了六个运动任务子类别的跨会话性能表现，其中(a)和(d)分别对应三类手部抓握任务的运动执行MM-G和运动想象MI-G，(b)和(e)分别对应六类手臂伸展任务的运动执行MM-R和运动想象MI-R，(c)和(f)分别对应两类腕部扭转任务的运动执行MM-T和运动想象MI-T。实验对比了传统基于CSP的方法（CSP+LDA和CSP+rLDA的平均性能）、现有监督学习模型（FBCNet、EEGNet、EEGViT和EEGConformer的平均性能）以及所提出的DMAE-EEG方法。


\begin{figure*}[htbp]
\centering
\includegraphics[width=\textwidth]{C5-8crosssess-CN.pdf}
\caption{跨会话泛化性能随校准样本数量变化的趋势}
\caption*{\centering (a) 三类手部抓握执行任务MM-G；(b) 六类手臂伸展执行任务MM-R；(c) 两类腕部扭转执行任务MM-T；(d) 三类手部抓握想象任务MI-G；(e) 六类手臂伸展想象任务MI-R；(f) 两类腕部扭转想象任务MI-T。}
\label{fig:5-8}
\end{figure*}

实验结果显示，DMAE-EEG在极少量自适应校准数据（仅需125秒，相当于每类50个样本）的支持下，即可在平衡准确率和Kappa系数两项关键指标上实现显著性能提升。在分类难度相对较低的三类任务中，DMAE-EEG在MM-G和MI-G任务上分别达到75.62\%和75.52\%的准确率，相比最优监督学习方法分别提升1.45\%和1.50\%。在更具挑战性的六类细粒度分类任务中，DMAE-EEG在MM-R和MI-R任务上的准确率分别为61.90\%和61.67\%，相比监督方法分别获得1.81\%和1.68\%的显著提升。在二分类腕部扭转任务中，DMAE-EEG表现尤为突出，MM-T和MI-T任务准确率分别达到84.44\%和80.14\%，相比监督方法分别提升3.20\%和1.34\%，已接近会话内最优性能水平。从Kappa系数角度分析，DMAE-EEG在所有任务中均表现出一致的优势，特别是在两类任务中表现最为显著，MM-T和MI-T的Kappa值相比监督方法分别提升10.47\%和8.65\%。这一结果表明DMAE-EEG编码器具备了学习稳定且可迁移特征表征的核心能力。此外，随着自适应校准样本数量的逐步增加，所有方法的性能均呈现稳步提升趋势，但DMAE-EEG表现出更优的样本效率和更快的收敛速度，充分验证了其在有限标注数据条件下的卓越适应能力。

从神经机制角度分析，跨会话泛化性能的显著提升主要归因于DMAE-EEG在大规模无标注预训练阶段所学习到的与会话无关的神经活动本质模式。这些深层特征表征能够有效捕获运动皮层区域的共性激活模式，而非依赖于特定会话的表面统计特性，从而在面对信号漂移、电极位置变化等现实挑战时仍能维持稳定的解码性能。结合少样本自适应校准机制，模型能够通过有限的目标域数据快速调整决策边界，有效弥合训练分布与测试分布之间的域间隙，为BCI系统的长期稳定应用提供了重要的技术保障。

\subsection{跨被试泛化能力的差异适应性} 

实际BCI系统部署场景中，设备往往需要在计算资源受限的边缘环境下为新用户提供服务，而无法针对每位用户进行大规模模型重训练。这要求BCI解码器具备强大的跨受试者泛化能力，能够快速适应个体间神经信号的固有差异。为了系统性评估DMAE-EEG编码器在多用户场景下的适应性和可扩展性，本研究基于MultiM11数据集设计了跨受试者评估实验。实验采用受试者分割策略，随机选择5名未参与预训练过程的受试者数据作为微调基准集，剩余5名受试者数据构成独立测试集，并在目标用户测试阶段通过10-50个样本片段（每类2-10个样本，总计25-125秒）进行个性化自适应校准。
图~\ref{fig:5-9}展示了跨受试者场景下的性能表现。实验系统对比了传统信号处理方法、现有深度学习监督模型以及所提出DMAE-EEG方法在不同校准样本规模下的性能表现。


\begin{figure*}[!ht]
\centering
\includegraphics[width=\textwidth]{C5-9crosssub-CN.pdf}
\caption{跨被试泛化性能随校准样本数量变化的趋势}
\caption*{\centering (a) 三类手部抓握执行任务MM-G；(b) 六类手臂伸展执行任务MM-R；(c) 两类腕部扭转执行任务MM-T；(d) 三类手部抓握想象任务MI-G；(e) 六类手臂伸展想象任务MI-R；(f) 两类腕部扭转想象任务MI-T。}
\label{fig:5-9}
\end{figure*}

实验结果表明，DMAE-EEG在跨受试者场景下展现出卓越的泛化能力和样本效率。在充分校准条件下（50个样本），DMAE-EEG在三类任务MM-G和MI-G上分别达到75.15\%和75.48\%的平衡准确率，相比最优监督学习方法分别提升1.62\%和0.91\%。在更具挑战性的六类细粒度分类任务中，MM-R和MI-R的准确率分别为61.55\%和60.95\%，相比监督方法分别获得1.10\%和2.06\%的显著提升。在二分类任务中，DMAE-EEG表现尤为突出，MM-T和MI-T任务准确率分别达到81.94\%和78.47\%，相比监督方法分别提升3.43\%和2.60\%。
更为重要的是，DMAE-EEG在极少样本校准条件下（仅10个样本，相当于25秒数据）仍能维持稳定的解码性能，体现出优异的样本效率。从Kappa系数分析可以看出，DMAE-EEG在所有任务中均获得一致性改善，特别是在两类任务中表现最为显著。随着个性化校准样本数量的逐步增加，DMAE-EEG展现出平滑且稳定的性能提升曲线，表明其能够有效利用增量标注数据进行渐进式适应，而非简单的过拟合现象。

从神经科学角度分析，跨受试者泛化能力的核心在于提取与个体差异无关的运动意图共性特征。尽管不同受试者在EEG信号的绝对幅值、频谱功率分布、空间拓扑模式等方面存在显著的个体差异，但运动皮层在执行特定运动任务时的相对激活模式具有跨个体的一致性。DMAE-EEG通过大规模无监督预训练学习到的深层表征能够有效捕获这些本质性的神经活动模式，从而为快速个体适应提供了坚实的特征基础。结合轻量级微调机制，模型能够在保持通用特征的同时，通过少量个性化数据调整决策边界以适应特定用户的信号特性，为实际BCI系统的大规模部署提供了可行的技术路径。

\subsection{执行与想象任务间的跨任务迁移}
运动执行与运动想象任务虽然在神经信号强度、时空分布模式以及伪迹水平方面存在显著差异，但两者在运动皮层、前运动皮层和顶叶皮层等关键脑区的激活模式具有相似性。运动执行过程中的EEG信号通常表现出更高的信噪比和更稳定的时间一致性，因为实际动作产生的运动相关电位具有更强的幅值和更清晰的时频特征。相比之下，运动想象信号往往受到更多内在噪声和认知负荷的影响，导致特征模式相对模糊。基于这一观察，本研究探索了利用高质量MM数据增强MI任务解码性能的可能性。

实验采用跨任务迁移学习策略，首先使用DMAE-EEG预训练编码器在MM数据上进行任务特定微调，随后通过少量MI样本进行目标域自适应校准。图~\ref{fig:5-10}展示了从MM任务向对应MI任务迁移的性能表现，其中(a)对应三类手部抓握任务(MM-G$\rightarrow$MI-G)，(b)对应六类手臂伸展任务(MM-R$\rightarrow$MI-R)，(c)对应两类腕部扭转任务(MM-T$\rightarrow$MI-T)。实验结果显示，当MI适应样本增加至每类50个时，DMAE-EEG在三个迁移任务上分别达到76.46\%、60.24\%和82.09\%的准确率，相比传统方法获得显著性能提升。特别是在手部抓握想象任务MI-G中，该方法展现出突出的实用价值，因为抓握想象作为自然直观的交互范式，在BCI控制的机械臂系统中具有重要应用前景，例如想象抓取特定物体可直接触发机械臂执行相应的抓握动作。


\begin{figure*}[htbp]
\centering
\includegraphics[width=1.0\textwidth]{C5-10crosstask-CN.pdf}
\caption{跨任务迁移性能随校准样本数量变化的趋势}
\caption*{\centering (a) 三类手部抓握任务MM-G$\rightarrow$MI-G；(b) 六类手臂伸展任务MM-R$\rightarrow$MI-R；(c) 两类腕部扭转任务MM-T$\rightarrow$MI-T。}
\label{fig:5-10}
\end{figure*}

图~\ref{fig:5-11}通过注意力可视化分析揭示了跨模态迁移学习的内在机制。可视化结果显示，仅在MM-G数据上微调的编码器所关注的脑区分布与完全在MI-G数据上训练的模型表现出显著的相似性，这表明两种任务确实激活了相似的神经网络。随着MI域适应过程的进行，模型的注意力模式逐渐向MI特异性区域聚焦，体现了预训练表征的良好可迁移性和适应性。

\begin{figure*}[htbp]
\centering
\includegraphics[width=\textwidth]{C5-11vis-CN.pdf}
\caption{跨任务迁移学习中的注意力分布变化}
\label{fig:5-11}
\end{figure*}

此外，DMAE-EEG展现出零样本跨任务识别能力。如图~\ref{fig:5-12}所示，在完全没有目标域标注数据的情况下，模型在所有六个跨任务中均取得了超越随机水平的性能。在相对简单的两类任务中，MM-T$\rightarrow$MI-T和MI-T$\rightarrow$MM-T分别达到59.67\%和59.75\%的准确率，相比50\%的随机基线获得近10个百分点的提升。在更具挑战性的多类任务中，三类任务(MM-G$\rightarrow$MI-G为39.76\%，MI-G$\rightarrow$MM-G为41.25\%)和六类任务(MM-R$\rightarrow$MI-R为23.14\%，MI-R$\rightarrow$MM-R为21.89\%)同样显著超越相应的随机基线33.3\%和16.7\%。这些结果充分验证了DMAE-EEG在提取跨任务共享特征方面的卓越能力。

\begin{figure*}[htbp]
\centering
\includegraphics[width=0.99\textwidth]{C5-12zeroshot.pdf}
\caption{零样本跨任务识别性能表现}
\caption*{\centering 注：横虚线表示对应任务的随机基线水平}
\label{fig:5-12}
\end{figure*}


零样本跨模态迁移能力的实现源于DMAE-EEG对运动神经网络内在组织结构的深层建模。在掩码自编码预训练过程中，模型被迫从部分观测的EEG信号中重构完整的时空模式，这一约束机制促使编码器学习到运动皮层激活的核心不变量。具体而言，无论是实际执行手部抓握动作还是想象该动作，初级运动皮层的对侧激活模式、辅助运动区的运动准备信号以及感觉运动皮层的节律抑制现象都表现出一定的一致性。DMAE-EEG通过自监督重构任务自然地捕获了这些跨模态共享的神经标记，而非依赖于任务特定的表面统计特征。这种抽象特征表征的有效性在不同复杂度任务中表现出差异化的迁移效果。在两类腕部扭转任务中，运动想象与执行间的神经活动模式相对简单且区分度较高，主要涉及对侧运动皮层的激活差异，因此零样本迁移性能接近60\%。相比之下，六类手臂伸展任务需要模型区分更细微的运动方向差异，这些差异不仅体现在空间拓扑上，还涉及不同肌肉群协调模式的神经编码，使得跨模态迁移面临更大挑战。值得注意的是，即使在最复杂的六类任务中，DMAE-EEG仍能维持超越随机基线约6-7\%的性能，表明模型确实学习到了有意义的跨模态特征映射关系。

\section{本章小结}
本章提出了基于去噪掩码自编码器的EEG表征学习方法DMAE-EEG，通过构建统一的自监督学习框架解决EEG分析中标注数据稀缺和噪声伪迹严重的核心问题。该方法整合神经科学先验知识与深度学习技术，设计了脑区拓扑异质划分策略显式编码跨区域拓扑差异并提升跨数据集对齐能力。划分方法有效处理了不同采集设备间电极配置和布局的差异性，为大规模多源数据的统一建模奠定了基础。去噪伪标签生成器通过重建型去噪任务生成稳健的教师信号，缓解噪声和伪迹对预训练目标的干扰，引导模型学习具有任务无关性和可迁移性的表征。这种设计有效区分了真实神经活动与各种干扰成分，显著提升了表征的鲁棒性和判别性。
在14个公开数据集上进行的大规模预训练和下游任务评估充分验证了DMAE-EEG的有效性和优越性。实验结果表明，该方法在EEG信号重建任务中显著降低了重建误差，信号质量得到明显提升。在运动意图识别任务中，DMAE-EEG展现出卓越的泛化能力，在跨会话、跨受试者和跨任务场景下均实现了显著的性能提升。特别地，该方法通过最少的校准数据即可实现快速适应，大幅减少了频繁重新校准的负担，显著提升了BCI系统在长期使用场景中的可行性和实用性。这些优势使得DMAE-EEG在减少对任务特定标注依赖的同时，显著改善了EEG系统的可扩展性和实用性。

% 基于本章研究成果，DMAE-EEG为EEG分析领域的自监督学习提供了新的技术路径和理论基础。该方法有效解决了传统有监督学习在EEG分析中面临的数据稀缺和噪声干扰问题，为构建更加智能、鲁棒和通用的EEG分析系统奠定了重要基础。通过显著降低对任务特定标注的依赖，DMAE-EEG推动了EEG技术在BCI、临床诊断和神经科学研究等领域的广泛应用。未来工作将进一步探索细粒度任务适应、参数高效性和多模态集成等方向，以扩展应用范围并构建更加自适应和易获取的BCI解决方案，为神经工程技术的产业化应用和临床转化提供有力支撑。